\chapter{Preguntas difíciles y su respuesta}
\label{cap:preguntas}

\section{Sobre el proyecto y los datos}
\begin{description}[style=nextline,leftmargin=0pt]
\item[¿Inventaste algún dato?] No. Lo que no existe es hueco; lo estimado lleva \codigo{_est}; cada cifra se rastrea hasta su
  archivo oficial en \codigo{docs/trazabilidad.md}, y las pruebas recalculan las cifras de los documentos.
\item[¿Por qué aparece Claude?] Porque usé un asistente de programación con inteligencia artificial. Las decisiones son mías
  y están documentadas con mis palabras; puedo explicar cada método (capítulo~\ref{cap:construccion}).
\item[¿Por qué el sur y no Bacalar o Tulum?] Sargazo récord en la franja de Tulum a Xcalak, crisis de Tulum ($-31.3\,\%$ de
  visitas en enero–julio de 2026) y la laguna de Bacalar en deterioro. El sur tiene capacidad ociosa: el 12.3\,\% de la
  población y el 1.4\,\% de los pasajeros de avión.
\item[¿Por qué tantos archivos?] Una pieza por archivo, para probar y explicar cada una por separado; los notebooks narran y
  usan esas piezas.
\item[¿De verdad necesitabas Spark?] Pandas podría con casi todo. Se usó por la arquitectura que pide la rúbrica, porque el
  DENUE nacional y el clima por hora ya son pesados, y porque el mismo código escala a un clúster.
\item[¿Qué haces si dos fuentes no coinciden?] Una fuente por lugar en toda su historia, la diferencia se declara, y se
  revisa en qué sentido va (aquí, conservador).
\end{description}

\section{Sobre los métodos}
\begin{description}[style=nextline,leftmargin=0pt]
\item[¿Por qué no promediaste porcentajes de ocupación?] Porque los meses con pocos cuartos pesarían igual que los grandes;
  se suman cuartos y luego se divide.
\item[¿Por qué pesos iguales en el índice?] No hay dato que diga qué componente pesa más, y la sensibilidad muestra que el
  estado cambia en 6\,\% de los casos como máximo con pesos de 0.5 a 1.5.
\item[¿Por qué validaste con origen móvil y no 80/20 al azar?] Porque en series de tiempo revolver al azar deja que el
  modelo vea el futuro; el origen móvil solo usa el pasado, como en la vida real.
\item[¿Tu clasificador sirve?] Le gana a ``igual que este mes'' por 4 aciertos de 156 y anticipa 8 de 30 cambios. En el sur
  casi no hubo cambios con qué probarlo, y así lo digo.
\item[¿Por qué no ARIMA o Prophet?] Las series tienen huecos largos y tramos cortos desde las reaperturas; esos métodos
  necesitan series continuas largas.
\item[¿Qué garantiza el rango del 90\,\%?] Si los errores futuros se parecen a los pasados, el dato cae dentro al menos 90
  de cada 100 veces. Medí la cobertura real: 80–94\,\%; la Ruta falló en 2023 y lo declaro.
\item[¿Por qué Poisson para las tormentas?] Son eventos raros e independientes en un periodo: es la distribución natural
  para contarlos.
\item[¿Por qué copias errores de un mismo origen en el Monte Carlo?] Porque los meses flojos vienen juntos; sortear cada
  mes por separado haría los años demasiado parejos.
\item[¿Qué es el golpe de tormenta de 50\,\%?] Un supuesto: solo hubo 3 meses con tormenta para medirlo; por eso se prueba
  con 0, 25 y 50\,\%.
\item[¿Por qué Isolation Forest?] Detecta lo raro sin que alguien defina ``raro''. Lo validé: marcó las 3 semanas de
  tormenta sin saberlo.
\item[¿Por qué léxico y no LDA en las reseñas?] Porque cada tema se define y se defiende palabra por palabra.
\end{description}

\section{Sobre la optimización y la campaña}
\begin{description}[style=nextline,leftmargin=0pt]
\item[¿Qué es lo estocástico de tu modelo de optimización?] La segunda etapa: la pausa depende del escenario, y el objetivo
  es el valor esperado sobre los escenarios.
\item[¿Qué significa el precio sombra?] Cuánto mejora el objetivo si se afloja una restricción una unidad. El del
  presupuesto es 1.59 visitantes por cada \$1,000.
\item[¿Cuánto cuesta la regla ambiental?] A este presupuesto, nada. Lo único que cuesta es el tope por canal (29.5\,\%).
\item[¿Por qué 70\,\% en Facebook?] Rinde 6.61 visitantes por cada \$1,000 contra 1.59 de Google; sin tope se iría el 100\,\%.
\item[¿De dónde sale la capacidad?] Es la capacidad probada: el mes más alto que cada lugar ya recibió; no es la capacidad
  física oficial.
\item[¿De dónde salen las personas?] INAH (nacional/extranjero), Unidad de Política Migratoria (país, sexo, mes) y ENDUTIH;
  los huecos se declaran.
\item[¿Qué pasa si la campaña funciona demasiado bien?] La capacidad y la temporada alta son restricciones del modelo, y la
  Torre pausa cada semana. Además, se puede prometer no pasar del 40\,\% de la capacidad sin perder visitantes.
\item[¿Probaron la página con personas?] No; se declara. Se hizo una auditoría automática con axe-core: 0 fallas en 4 vistas.
\end{description}
